\chapter{代数式}\label{ch:g7:literal}

运算中的字母表示一个尚不知道的数，或
同时表示\emph{所有}的数：写成 $P = 4c$，
即可用一个式子给出每个正方形的周长。本章讲字母的书写约定、
代入数值的方法，以及简单代数式的化简，
这是代数的开端，
将在\cref{ch:g8:equations}中继续。

\section{用字母书写}

\begin{definition}[代数式]\label{def:g7:literal:def}
称\emph{代数式}\index{代数式}为一种运算，
含有一个或多个字母，每个字母表示一个数。
书写约定如下：
\begin{itemize}
  \item 在字母或括号前，乘号 $\times$ 可以省略：
        $3 \times a = 3a$, \ $a \times b = ab$, \
        $2 \times (x + 5) = 2(x + 5)$;
  \item 但两个数之间的乘号须保留：$3 \times 5$ 不能
        写成 $35$！
  \item $a \times a$ 写作 $a^2$（“$a$ 的平方”），
        $a \times a \times a = a^3$;
  \item $1a$ 写作 $a$，$0a$ 则为 $0$。
\end{itemize}
\end{definition}

\begin{example}\label{ex:g7:literal:write}
按约定改写下列式子：
\[
5 \times x + 3 \times y = 5x + 3y,
\qquad
a \times 7 = 7a \ \text{（数字在前）},
\qquad
x \times x \times 4 = 4x^2 .
\]
\end{example}

\begin{example}[公式也是代数式]\label{ex:g7:literal:formulas}
对于长为 $L$、宽为 $w$ 的长方形：
周长 $P = 2(L + w)$，面积 $A = Lw$。一个式子概括了
无穷多个长方形，这正是字母的
作用所在。
\end{example}

\begin{omfigure}
\begin{tikzpicture}[scale=0.8]
  \draw[thick] (0,0) rectangle (4.4,2.0);
  \node[below, font=\small] at (2.2,0) {$L$};
  \node[left, font=\small] at (0,1.0) {$w$};
  \node[font=\small] at (2.2,1.0) {$A = Lw$};
  \node[above, font=\small, omDef] at (2.2,2.05) {$P = 2(L + w)$};
\end{tikzpicture}

{\small 两个式子对每个长方形都成立：代入
$L$ 和 $w$ 的值，即可得到数值。}
\end{omfigure}

\section{代入求值}

\begin{method}[代数式求值]\label{met:g7:literal:substitute}
给定字母的值，求代数式的值，步骤如下：
\begin{enumerate}
  \item 将代数式中省略的乘号 $\times$ 补回来：
        $3a + a^2$ 改写为 $3 \times a + a \times a$；
  \item 把每个字母换成它的取值，
        若有需要，则保留括号；
  \item 按照
        \cref{ch:g7:priorities}中的运算顺序进行计算。
\end{enumerate}
\end{method}

\begin{example}\label{ex:g7:literal:substitute}
求 $E = 3x + x^2$ 在 $x = 5$ 时的值：
\[
E = 3 \times 5 + 5 \times 5 = 15 + 25 = 40 .
\]
求 $F = 2(a + 3b)$ 在 $a = 4$、$b = 0.5$ 时的值：
\[
F = 2 \times (4 + 3 \times 0.5) = 2 \times (4 + 1.5) = 2 \times 5.5
= 11 .
\]
\end{example}

\section{化简代数式}

\begin{proposition}[合并同类项]\label{prop:g7:literal:reduce}
含有相同字母部分的项，可以把它们的
系数相加合并（这就是分配律，
由\cref{thm:g7:priorities:distributivity}反向得到）：
\[
5x + 3x = (5 + 3)x = 8x,
\qquad
7a - 2a = 5a .
\]
字母部分不同的项（如 $3x$ 与 $5y$，或 $x$ 与
$x^2$）不能合并。
\end{proposition}

\begin{example}\label{ex:g7:literal:reduce}
化简 $E = 4x + 7 + 3x + 2$，并将同类项归在一起：
\[
E = (4x + 3x) + (7 + 2) = 7x + 9 .
\]
注意：$7x + 9$ \emph{不}是 $16x$！取 $x = 1$ 检验：
$4 + 7 + 3 + 2 = 16$ 与 $7 \times 1 + 9 = 16$ 相等，但 $16x$ 也
等于 $16$。取 $x = 2$ 检验：原式为
$8 + 7 + 6 + 2 = 23$，$7 \times 2 + 9 = 23$，而
$16 \times 2 = 32$。化简结果 $7x + 9$ 正确，$16x$ 错误。
\end{example}

\begin{method}[检验等式]\label{met:g7:literal:test}
要检验两个代数式对每个取值是否可能相等：
\begin{enumerate}
  \item 在两个代数式中代入相同的值，至少两次（例如\
        $x = 2$ 与 $x = 10$；应避免取 $0$ 和 $1$，它们会掩盖太多
        差异）；
  \item 只要结果哪怕有一次不同，两个代数式就
        \emph{不}相等 --- 一个反例就足以推翻一般性的论断；
  \item 若结果一致，则该等式仅是\emph{可能成立}的：
        证明则需借助代数运算（展开、化简）。
\end{enumerate}
\end{method}

\begin{example}\label{ex:g7:literal:test}
$2(x + 3)$ 与 $2x + 6$ 对每个 $x$ 是否相等？展开即可证明：
$2(x+3) = 2x + 2 \times 3 = 2x + 6$。$(x + 1)^2$ 是否与
$x^2 + 1$ 相等？取 $x = 2$ 检验：$(2+1)^2 = 9$，而 $2^2 + 1 = 5$。答案为否 ---
一个反例即可完整证明“否”。
\end{example}

\section{建立代数式}

\begin{example}[从情境到代数式]\label{ex:g7:literal:produce}
一辆出租车收取 $4$ 欧元的起步价，另加每千米 $2$ 欧元。若行程为
$k$ 千米，则车费为
\[
p = 4 + 2k \quad\text{欧元}.
\]
当 $k = 7$ 时，$p = 4 + 14 = 18$ 欧元。该代数式\emph{就是}
一般性的答案；代入不同的数值即可得到每一个具体结果。
\end{example}

\begin{example}[数字魔术]\label{ex:g7:literal:trick}
“想一个数，把它加倍，加上 $10$，再除以 $2$，减去
你所想的那个数。”设这个数为 $n$，按下列步骤计算：
\[
n \ \to\ 2n \ \to\ 2n + 10 \ \to\ \frac{2n + 10}{2} = n + 5 \ \to\
n + 5 - n = 5 .
\]
人人得到的结果都是 $5$，无论 $n$ 原为何数 ---
字母道出了这个“魔术”的奥秘。
\end{example}

\section{练习}

\begin{exercise}[$\star$]\label{exo:g7:literal:1}
把下列式子写得更简洁：
\[
7 \times x, \qquad
y \times 3, \qquad
a \times b \times 5, \qquad
x \times x, \qquad
2 \times (a + 4).
\]
\end{exercise}

\begin{exercise}[$\star$]\label{exo:g7:literal:2}
把 $\times$ 补写回去：
\[
5ab, \qquad
3(x + 2), \qquad
x^2 + 4x, \qquad
2a^3 .
\]
\end{exercise}

\begin{exercise}[$\star$]\label{exo:g7:literal:3}
当 $x = 3$ 时，求下列各式的值：\ $5x + 1$；\quad $x^2$；\quad $2x^2 - x$；
\quad $10 - 2x$。
\end{exercise}

\begin{exercise}[$\star$]\label{exo:g7:literal:4}
求 $3a + 2b$ 与 $a b + 5$ 在 $a = 4$，$b = 1.5$ 时的值。
\end{exercise}

\begin{exercise}[$\star$]\label{exo:g7:literal:5}
化简：
\[
8x + 5x, \qquad
9a - 4a + a, \qquad
6y + 2 + 3y + 7, \qquad
5t + 3 - 2t - 3 .
\]
\end{exercise}

\begin{exercise}[$\star$]\label{exo:g7:literal:6}
写出下列各量的公式：边长为 $c$ 的正方形的周长；
底边为 $b$、腰长为 $s$ 的等腰三角形的周长；
$n$ 个面包的价格，每个 $1.20$ 欧元。
\end{exercise}

\begin{exercise}[$\star$]\label{exo:g7:literal:7}
将每个代数式与一种描述配对：$2n$；$n + 2$；$n^2$；
$\frac n2$。描述为：“$n$ 的一半”；“$n$ 的两倍”；
“$n$ 的平方”；“比 $n$ 多 2”。
\end{exercise}

\begin{exercise}[$\star\star$]\label{exo:g7:literal:8}
利用 \cref{met:g7:literal:test}，判断下列等式是否
对每个 $x$ 成立，并给出证明或反例：
\[
3(x + 4) = 3x + 12; \qquad
2x + 5 = 7x; \qquad
x + x = x^2 .
\]
\end{exercise}

\begin{exercise}[$\star\star$]\label{exo:g7:literal:9}
一个长方形的宽为 $w$，其长比
宽多 $3$ cm。
\begin{enumerate}
  \item 用 $w$ 分别表示长与周长，
        并化简。
  \item 当 $w = 5$ cm 时，分别直接计算周长，
        并用化简后的公式计算。结果是否相同？
\end{enumerate}
\end{exercise}

\begin{exercise}[$\star\star$]\label{exo:g7:literal:10}
用字母重做\cref{ex:g7:literal:trick}中的小把戏：“想一个数，
加上 $6$，乘以 $3$，减去 $18$，再除以
你开始时所想的数”（假设该数不为零）。每个人算出的结果是什么？
\end{exercise}

\begin{exercise}[$\star\star\star$]\label{exo:g7:literal:11}
下图是边长为 $a$ 的正方形，去掉一个边长为 $b$ 的小正方形角
（$b$ 小于 $a$）。写出其面积的两种不同表达式：一种用差表示，
另一种将 L 形剪成两个长方形，
并检验当 $a = 5$，$b = 2$ 时两者相等。
\end{exercise}

\section{问题：火柴神谕}

\begin{problem}[{周末问题 --- 增长的图形规律及其第
$n$ 项：同一规律的三个公式，以及著名的
在第六步失效的规律}]
\label{pb:g7:literal:1}
用火柴摆成一行正方形：先摆一个正方形，再摆两个，再摆
三个\,\dots\ 第一百个图形需要多少根火柴？
一根一根地数会花掉整个周末；而\emph{
关于 $n$ 的公式}能立即给出所有图形的结果，
这正是字母的用途
（见\cref{ex:g7:literal:formulas}）。但要当心：本题以一个
著名规律收尾，它在前五步看似公平，
随后却辜负了每一个信任它的人。

\medskip
\noindent\textbf{第 I 部分 --- 一行正方形。}
图形 $1$ 是一个火柴正方形（$4$ 根火柴）；此后每个图形都
增加一个与前一图形共享一条边的正方形。
\begin{enumerate}
  \item 画出图形 $1$ 至 $4$，并数出各图形的火柴根数。
  \item 说明为什么每增加一个新正方形恰好需要 $3$ 根
        额外的火柴，并利用这一点预测图形
        $10$ 的火柴根数，不借助任何公式。
  \item 论证公式：图形 $n$ 需要 $3n + 1$ 根火柴。
        （那个孤零零的“$+1$”从何而来？）并
        将其与第 1 题的结果对照。
  \item 阿米尔的思路不同：“第一个正方形用 $4$ 根火柴，
        之后每个正方形用 $3$ 根，共 $n - 1$ 个”，得到
        $4 + 3(n - 1)$。展开并化简他的表达式
        （见\cref{prop:g7:literal:reduce}）：他的结果与
        第 3 题是否一致？
  \item 莉娜用第三种方法计数：“$2n$ 根水平火柴，以及
        $n + 1$ 根竖直火柴”。写出她的公式并化简，
        再得出结论。然后回答开头的问题：
        图形 $100$ 需要多少根火柴？
\end{enumerate}

\medskip
\noindent\textbf{第 II 部分 --- 三角形、桌子与
逆推。}
\begin{enumerate}[resume]
  \item 一行火柴三角形（每个新三角形都紧靠前一个）：数出
        $1$、$2$、$3$ 个三角形所用的火柴根数，
        然后论证公式 $2n + 1$，它给出
        $n$ 个三角形所用的火柴根数。
  \item 逆推：恰好用 $49$ 根火柴时，这一行最多能有多少个
        三角形？（请对公式逐步逆用，
        每次撤销一步。）
  \item 正方形宴会桌，每条空闲的边坐一位客人。当
        将 $n$ 张桌子拼成一长排时，说明座位数公式为何是
        $2n + 2$，并对 $1$、$2$、
        $3$ 张桌子的情况进行检验。
  \item 一场聚会预计有 $30$ 位客人。摆成一排需要多少张桌子？
        餐饮承办方则将同样的桌子分成
        \emph{两}排，每排七张：现在能坐下多少位客人？
        请用公式解释增加的人数。
  \item 点阵规律：图形 $n$ 是一个 L 形点阵，高 $n$ 个点，
        宽 $n$ 个点（拐角只计一次）。数出图形 $1$、$2$、$3$ 的点数，
        求出公式，并说出它所产生的数。
        （把这些 L 形叠放起来，可以铺成一个正方形——
        这一故事在
        \cref{pb:g8:equations:1}中继续。）
\end{enumerate}

\medskip
\noindent\textbf{第 III 部分 --- 信任、检验与证明。}
\begin{enumerate}[resume]
  \item 一位同学声称第 I 部分中图形 $n$ 需要
        $3(n + 2) - 5$ 根火柴。请通过展开并化简来判定这一说法，
        无需逐个检验。
  \item “取平方会使数变大，所以 $n^2 \geq n$ 对每个数
        $n$ 都成立。”分别对 $n = 2$、$3$、
        $10$ 进行检验，然后找出一个使它不成立的数。
        三次成功的检验证明了什么，而
        唯一的一次失败又证明了什么（见\cref{met:g7:literal:test}）？
  \item 在圆周上标出若干点，并画出它们之间的\emph{每一条}线段，
        然后数出圆内的区域个数。
        对 $2$、$3$、$4$、$5$ 个点分别进行（点的位置要参差不齐，
        使圆内任何三条线段都不交于同一点）。
        记录各次的数目。这一规律悄悄暗示了
        什么公式？
  \item 现在取 $6$ 个点，画得大一些、位置不均匀，
        格外小心地计数（区域很多，而且很小）。你数出了
        多少个区域？这一猜想
        后来怎样了？
  \item 用两句话概括其中的道理：要
        \emph{证明}倍增猜想需要什么，以及第 I、II 部分的火柴公式为何
        从未面临这样的危险？
        让我们用一个已被证明的预测来庆祝一下：
        一行 $2\,026$ 个正方形需要多少根火柴？
\end{enumerate}
\end{problem}
